%% =====================================================================
%%  T-17-04  The Second Self
%%  -------------------------------------------------------------------
%%  One idea per file. Core is mandatory. Slots in canonical order.
%%  Max 700 words. No layout commands. No filler. New source material
%%  for this entry goes in sources/B3/T-17-04.tex -- never by
%%  rewriting this file (docs/RULES.md R0.1).
%%  =====================================================================
%% @kind: topic
%% @id: T-17-04
%% @title: The Second Self
%% @subtitle: Re-create yourself --- be the master of your own image
%% @chapter: c17
%% @order: 40
%% @status: stable
%% @sources: B3
%% @updated: 2026-09-19

\topic{T-17-04}{The Second Self}{Re-create yourself --- be the master of your own image}

\begin{core}
Refuse the part the world handed you. Forge a fresh identity --- one that draws the eye and never wearies the room --- and own your image instead of renting it from other people's mouths; fold dramatic devices into your public gestures until the character reads larger than life. The self presented is a constructed artefact; the only question is who holds the tools.
\end{core}
\begin{mechanism}
B3's law treats self-presentation as authorship. The mechanism is role-taking: audiences assign you a part --- the competent one, the difficult one, the one who almost matters --- and then interact with the part, feeding back the treatment the part deserves. Accept the assigned role and the definition compounds; forge a new one and you are interacting with the audience's expectations rather than from inside them. The craft is theatrical in the source's own terms: invent dramatic devices for public gestures, control the stagecraft (T-17-05), and take the mask seriously as construction --- the reversal warns that bad theatre is still bad theatre, and even appearing natural requires art. The chapter-family connection: this is the affirmative parent of the concealment entries (T-13-03), the entry craft of the courtier (T-12-05), and the mobility of the final chapter (T-28-01) --- the second self is the instrument those disciplines play.
\end{mechanism}
\begin{conditions}
Strongest for public-facing roles --- leadership, brands, performance, any life lived before audiences --- and at identity breakpoints: new city, new job, new market, where no local definition of you exists yet to fight. It weakens among intimates, where the second self reads as distance, and in contexts that audit sincerity; the construct must be possible, or the audience that meets the reality will invoice the gap.
\end{conditions}
\begin{application}
  \item Write the role before the room writes it: three traits you want to be knowable for, and the behaviours that display each.
  \item Choose the breakpoint. Roles are cheapest to install where nobody holds a definition of you yet.
  \item Be inconsistent on purpose at the level of surface, consistent at the level of story --- predictability of character, unpredictability of behaviour (T-19-04).
  \item Rehearse the devices: the entrance, the pause, the signature gesture. Drama is a skill with scales, not a gift.
  \item Keep one audience who sees the workshop. Total performance has no editor, and unedited performers drift into parody.
\end{application}
\begin{feedback}
  \item Working: strangers treat you as the role within minutes --- the definition is installed and working for you.
  \item Working: the old role's expectations stop arriving. The re-write took.
  \item Failing: the role demands more than you budgeted --- the character is running your calendar. Authorship has inverted.
  \item Failing: audiences sense the seams and read inauthenticity. The construct outran the possible; shrink it to fit.
\end{feedback}
\begin{failure}
The second self has a well-documented failure family: the construct that detaches (the performer who can no longer locate the person), the gap invoice (an audience eventually meets the un-constructed self and prices the difference as fraud), and role exhaustion --- maintaining a larger-than-life character is a permanent tax on the attention the character exists to draw. The source's guard is the only guard: the art must be good enough that the art disappears.
\end{failure}
\begin{countermeasures}
  \item When someone seems suspiciously well-fitted to their role, assume authorship and ask when they chose it. The answer dates the construct.
  \item Audit the roles assigned to you --- by family, by office, by history. Whose handwriting is on your current part?
  \item Give performers their invoice honestly: judge the work, not the self, and say which you are judging.
  \item In your own re-creations, keep the gap between construct and capacity small enough to survive a surprise audit (T-09-02).
\end{countermeasures}

\sources{B3}
\seesources
\seealso{T-17-05, T-13-03, T-28-01}
